\documentclass[UTF8,a4paper,11pt,fontset=fandol]{ctexart}
\usepackage[margin=2.2cm]{geometry}
\usepackage{amsmath,amssymb,booktabs}
\setlength{\parskip}{0.3em}
\linespread{1.1}

\begin{document}
\pagestyle{plain}

\begin{center}
  {\Large\bfseries 第三次上机作业：例 4.15 的网页排序}
\end{center}

\section{原题抄录}

编程重现教材例 4.15：取 $\alpha=0.85$，
运用第 4.5.2 节的方法对图 4.7 中的六个节点进行网页排序。

\section{代码的结构、算法说明}

\subsection*{（1）构造转移概率矩阵}

沿用教材记号，$a_{ij}=1$ 表示节点 $v_j$ 指向节点 $v_i$。
图 4.7 的邻接矩阵为
\[
A=
\begin{pmatrix}
0&0&1&0&1&1\\
1&0&1&0&1&0\\
0&1&0&0&0&0\\
1&1&0&0&0&0\\
0&0&0&0&0&0\\
0&0&0&0&1&0
\end{pmatrix}.
\]
程序使用 \texttt{np.sum} 计算行和与列和：
\[
r_i=\sum_{j=1}^{n}a_{ij},\qquad
c_j=\sum_{i=1}^{n}a_{ij}.
\]
其中 $r_i$ 表示指向 $v_i$ 的节点数量，
$c_j$ 表示 $v_j$ 所指向的节点数量。

取 $n=6$、$\alpha=0.85$，逐列构造
\[
p_{ij}=
\begin{cases}
\displaystyle \frac{\alpha a_{ij}}{c_j}
+\frac{1-\alpha}{n},&c_j\ne0,\\[6pt]
\displaystyle \frac1n,&c_j=0.
\end{cases}
\]
节点 $v_4$ 无出链，因此 $P$ 的第 4 列全部取 $1/6$。
所得 $P$ 是严格为正的列随机矩阵，
对应的 Markov 链不可约、非周期，
因而存在唯一平稳分布，且极限分布收敛。

\subsection*{（2）幂法计算与排序}

取均匀初始分布
\[
\pi_0=\frac16(1,1,1,1,1,1)^{\mathrm T},
\]
通过幂法迭代
\[
\pi_t=P\pi_{t-1}.
\]
当相邻两次迭代满足
\[
\|\pi_t-\pi_{t-1}\|_\infty<10^{-12}
\]
时停止，最大迭代次数设为 10000。

程序使用 \texttt{np.argsort(-pi)}
将平稳概率从大到小排序，
并将 Python 下标加 1，得到教材中的节点编号。
函数 \texttt{print\_array} 统一输出向量和矩阵，
浮点数保留小数点后四位。

\section{输出结果呈现}

邻接矩阵的行和、列和分别为
\[
r=(3,3,1,2,0,1)^{\mathrm T},\qquad
c=(2,2,2,0,3,1)^{\mathrm T}.
\]

转移概率矩阵保留四位小数后为
\[
P\approx
\begin{pmatrix}
0.0250&0.0250&0.4500&0.1667&0.3083&0.8750\\
0.4500&0.0250&0.4500&0.1667&0.3083&0.0250\\
0.0250&0.4500&0.0250&0.1667&0.0250&0.0250\\
0.4500&0.4500&0.0250&0.1667&0.0250&0.0250\\
0.0250&0.0250&0.0250&0.1667&0.0250&0.0250\\
0.0250&0.0250&0.0250&0.1667&0.3083&0.0250
\end{pmatrix}.
\]

幂法得到的平稳分布为
\[
\pi\approx
(0.2124,0.2366,0.1612,0.2514,0.0606,0.0778)^{\mathrm T}.
\]
网页排名结果如下：
\begin{center}
\begin{tabular}{ccc}
\toprule
排名 & 节点 & 平稳概率\\
\midrule
1 & $v_4$ & 0.2514\\
2 & $v_2$ & 0.2366\\
3 & $v_1$ & 0.2124\\
4 & $v_3$ & 0.1612\\
5 & $v_6$ & 0.0778\\
6 & $v_5$ & 0.0606\\
\bottomrule
\end{tabular}
\end{center}

\section{对输出结果的分析和结论}

计算所得平稳分布在四位小数显示精度下与教材例 4.15 一致，
节点排名为
\[
v_4>v_2>v_1>v_3>v_6>v_5.
\]

节点 $v_1$ 和 $v_2$ 均有较多入链，排名较高；
二者又同时指向 $v_4$，使 $v_4$ 获得最高排名。
节点 $v_5$ 没有入链，排名最低。
这表明网页的重要性不仅取决于入链数量，
还与指向它的网页的重要性有关。

本实验按照教材构造转移概率矩阵，
通过幂法计算平稳分布并完成节点排序，
成功重现了例 4.15 的结果。

\end{document}